\title{Collaborative Proposal --- ITR/IIS: Cross-modality Heterogeneous
Web Source Integration for High Quality Internet Information Access}

\newcommand{\MEGA}{\mathsf{MEGA}}
\newcommand{\SVMActive}{\mathsf{SVM}_{\mbox{\scriptsize{\textit{Active}}}}}


\section{Overview and Significance}
\label{sec-overview}

In this project we intend to develop fundamental technologies that 
will enable the users to access information
seamlessly, intuitively and quickly from the Web. 
The advent and success  of the World Wide Web has demonstrated 
the enormous opportunity that an easy-to-use and always-available
information system can bring. The users can access any information
at any time without any geographical constraints.
However, the very success of the Web also makes its current limitation
painfully obvious:

\begin{enumerate}

\item {\em Keyword-only search}:
  Internet search engines provide a great help
  to the users when the users look for Web pages on a certain topic.
  Simply by typing a couple of keywords, the user can quickly
  find relevant pages. However, the current search engines
  only provide {\em keyword} inferfaces, so the user is helpless
  when her ``information need'' cannot be described 
  by mere keywords. The rich set of multimedia information 
  on the Web is being totally ignored.

\item {\em Difficult query articulation}:
  Again, given her information need, the {\em user} has to
  come up with the appropriate keywords (or queries) that 
  describe her need succinctly. This task if often too difficult
  especially for novice users, leading to the user's frustration.
  How can a user describe ``I want to find more about paintings
  similar to the one that I saw a couple of days ago?''

\item {\em Stale and untimely information}:
  The Web is inherently dynamic medium: New pages
  are constantly created, obsolete pages are continually deleted
  and existing pages are updated over time.
  Therefore, it is often difficult to access ``up-to-date'' information,
  unless the user knows exactly where the information is located.
  For example, most search engines refresh their
  indexes periodically -- typically less than once a month --
  so recently created Web pages are often missing in their indexes. 
  In order to access these pages,
  the user needs to go directly to the appropriate Web site.
  Also, many of the pages returned by search engines
  are often obsolete and do not exist on the Web.

\item {\em Growing amount of dynamic information}:
  More and more information on the Web is dynamically generated.
  Often, the user can access this dynamic information
  only by visiting individual Web site and submitting queries,
  because search engines typically skip this ``hidden Web.''
  The pages cannot be reached by following links,
  so it is difficult for the search engines discover
  these pages. Clearly, it is both time consuming and difficult
  for the user to visit every Web site and issue
  queries to each of them to collect all information.
  
\end{enumerate}

To overcome these shortcomings, our project will develop new
technologies that can {\em discover} and {\em integrate} 
new and high-quality information from the Web in a timely manner,
support ``cross-modality'' search on multimedia objects, and
{\em learn} user's ``query concepts'' actively and quickly.
Towards these goals, we will organize our research around three main 
thrusts. Each thrust addresses a subset of the issues,
and as a whole our technologies will make the World Wide Web
easier to access even by naive users.

\begin{enumerate}
\item {\em Intelligent query concept formulation through active learning}.
 In order to make query formulation easy and complete, it is both
 necessary (for capturing abstract and subjective concepts) and
 desirable (for alleviating difficulties from specifying complex query
 concepts) to build intelligent search engines that can quickly {\em
 learn} users' query concepts through {\em active learning}.  We plan
 to extend our active learning algorithms $\MEGA$ 
 and $\SVMActive$ \cite{LCL01,SVM01},
 and develop new algorithms for
 cross-modality query-concept learning.

\item {\em Support of cross-modality search}. Depending on
  the domain of the information being sought for, the user
  should be able to query a database (or a search engine) 
  in different ``modes.''
  For example, if the user is looking for paintings similar
  to a particular one at hand, the colors, textures and ambience of images
  might be the best way to ``capture'' the user's ``query concept.''
  On the contrary, if the user is looking for Impressionism paintings,
  the keyword ``impressionism'' might be the best way to capture the user's
  concept. 

  We will develop technologies that can provide multiple ``modes''
  of interaction to the user. In particular, the user will
  be able to ``switch back and forth'' between different modes of
  interaction according to her varying need. 
  We will study techniques that can automatically extract
  textual meta data for multimedia objects from the Web.

\item {\em Efficient data collection and refresh}.
  We will develop algorithms for discovering ``important''
  and ``relevant'' information quickly from the Web
  and for downloading this information efficiently.
  We will also study mechanisms to track
  changes on the Web and refresh our database efficiently.
  We will also study how to {\em integrate} dynamic Web information
  seamlessly, so that the users can access the information 
  at a {\em single} point through a uniform interface.
\end{enumerate}

Our three research thrusts are both complementary and synergistic: Our
cross-modality paradigm will enrich Web data representation.  The
active learning algorithms will enable users to articulate their
``information need'' more intuitively and effectively, and thereby
make query formulation more personalized and complete.  And a timely
and comprehensive data collected from the Web will enhance the query
results.  The faculty involved in this proposal will all work on more
than one thrust.
%We have a record of sharing research interests, working together
%and joint publications. 
Our thrusts are a convenient way to organize our research agenda, but
are certainly not meant to partition the project.


\subsection{Example scenario}

{\bf (Now, our example scenario here. Junghoo 11/1)}

Consider that CIA wants to identify all information sources 
containing atomic bomb design in the Internet.
(This example was provided by In-Q-Tel, a CIA branch at Palo Alto.)
Unlike that most queries can be satisfied by just a couple of
good results, queries of this kind demand high recall.
To achieve high recall, a search engine
faces three hurdles:
comprehending the query concept completely, 
characterizing the data at hand adequately, 
and integrating with as much useful sources in the Internet 
as possible.
These three hurdles for achieving the performance goal are
precisely caused by the the three problems enumerated above.
This project aims to investigate novel 
techniques in three thrusts to overcome these three hurdles.



\subsection{Impact}
{\bf (We may need to revise this section, too. Junghoo 11/1)}

This project could have substantial impact upon image retrieval in
general, and upon education and military intelligence in particular.
For improving education, rapid proliferation of multimedia content in
digital libraries and on the Web underscores the increasing importance
of having effective cross-modality search tools.  Intelligent
query-concept learners will directly or indirectly make traditional
text-based information retrieval easier and more personalized.
Directly, a text collection can employ an intelligent learner to
better capture users' query concepts.  Indirectly, for instance,
multimedia data can be added to a text collection so that searches can
be conducted through interfaces that contain pictures and graphics. As
a result, even young students who have not learned Boolean algebra can
use images and graphics to search for stories and books.  For military
intelligence, the present world situation makes the gathering of
intelligence more essential than ever.  The ability to find
information that fit a query profile in a massive database offers a
strategic advantage to those who can assemble such information quickly
and accurately.



%%% Local Variables: 
%%% mode: latex
%%% TeX-master: "main"
%%% End: 
